% TEMPLATE PARA ESPECIALIZAÇÃO
% Instututo Tecnológico Vale (ITV)
% Ouro Preto - 2023
% Idealizado por Tatianna Beneteli

% Baseado nos projetos: 
% http://tcc.tsi.gp.utfpr.edu.br/paginas/modelos-latex-da-utfpr
% De: Diego Marczal e Michael Vornes
% Template para TCC - Universidade Federal do Pará
% De: Sérgio Custódio e Kelvin Pinheiro

\documentclass[oneside]{itv}
% oneside = Impresssão apenas na frente
% twoside = Impressão em frente e verso

%-----------------------------------------------------------------------
% PREENCHER OS CAMPOS ABAIXO:
%-----------------------------------------------------------------------

\titulo{Título do Trabalho}
\subtitulo{Subtítulo (se houver)} % Se não houver, deixar comentado (%)

\autor{Nome completo do autor} 

\local{Ouro Preto, MG} % separar por vírgula, caixa baixa
\data{09}{2023}
\diadefesa{02}

\projeto{Projeto / Relatório dos resultados parciais / Trabalho de conclusão de curso}
% Preencher com uma das opções: Projeto / Relatório dos resultados parciais / Trabalho de conclusão de curso

\tituloEspecializacao{Automação para Processos de Mineração}
% Exemplos:
% Lavra de Minas a Céu Aberto
% Beneficiamento Mineral
% Engenharia de Geotecnia Avançada
% Desgaste, Atrito e Lubrificação
% Automação para Processos de Mineração
% Sistemas Inteligentes: Ênfase em Ciência de Dados

\areaconcentracao{Otimização de Rotas}

% Dados do Orientador
\orientador{Prof. Nome completo do orientador} % Para orientadora, usar opção abaixo
%\orientador[Orientadora:]{Profa. Nome completo da orientadora}
\tituloabrevOrientador{D.Sc.}
\tituloOrientador{Titulação}
\instOrientador{Instituto Tecnológico Vale (ITV)}

% Dados do Coorientador
% Se não houver, será necessário comentar o trecho no arquivo "itv.cls" (linhas 329-331), além do trecho abaixo
\coorientador{Prof. Nome completo do coorientador} % Para coorientadora, usar opção abaixo
%\coorientador[Coorientadora:]{Profa. Nome completo da coorientadora}
\tituloabrevCoorientador{D.Sc.}
\tituloCoorientador{Titulação}
\instCoorientador{Instituição do coorientador (SIGLA)}

% Dados da banca
\nomePrimeiromembro{Prof. Nome do Avaliador interno}
\tituloPrimeiromembro{Titulação}
\instPrimeiromembro{Instituição (SIGLA)}
 
\nomeSegundomembro{Prof. Nome do Avaliador externo}
\tituloSegundomembro{Titulação}
\instSegundomembro{Instituição (SIGLA)}

% Para incluir um terceiro membro na banca, será necessário "descomentar" o trecho no arquivo "itv.cls" (linhas 341-343)
%\nomeTerceiromembro{Prof. Nome do terceiro Avaliador (se houver)}
%\tituloTerceiromembro{Titulação}
%\instTerceiromembro{Instituição (SIGLA)}

%-----------------------------------------------------------------------
% ESTRUTURA DO DOCUMENTO (cuidado ao realizar modificações)
%-----------------------------------------------------------------------

\begin{document}
	\pretextual
	\imprimircapa
    \imprimirfolhaderosto{}
    %\includepdf{fichacatalográfica.pdf}
    
    \imprimirfolhadeaprovacao{}
    % Este comando insere uma Folha de aprovação em branco, elemento obrigatório da NBR 14724/2011. Você pode utilizar este modelo até a aprovação do trabalho. Após isso, substitua todo o conteúdo deste arquivo por uma imagem da página assinada pela banca com o comando abaixo:
    %\includepdf{folhadeaprovacao_final.pdf}

    \include{Textos/pre-textuais/dedicatoria} % opcional
	\include{Textos/pre-textuais/agradecimentos}
	\include{Textos/pre-textuais/epigrafe} % opcional

    \OnehalfSpacing
    \include{Textos/pre-textuais/resumo}
	\include{Textos/pre-textuais/abstract}
	
    \listoffigures*
    \newpage
    \listoftables*
    %\listofalgorithms
    \include{Textos/pre-textuais/lista-siglas}
    \tableofcontents*
    \newpage
    
	\textual
	\include{Textos/textuais/1-Introducao}
	\include{Textos/textuais/2-ReferencialTeórico}
	\include{Textos/textuais/3-Metodologia}
	\include{Textos/textuais/4-ResultadosDiscussões}
	\include{Textos/textuais/5-Conclusões}
    \include{Textos/textuais/6-Sugestões}
    \include{Textos/textuais/7-Cronograma}

    % Orientações de uso do Template - Após leitura, a linha abaixo deve ser comentada
	\include{Textos/Orientacoes}
	
	\postextual
    \bibliography{referencias}
    \bibliographystyle{abntex2-alf}
	
	\include{Textos/pos-textuais/apendices} % Opcional (apêndices)
	\include{Textos/pos-textuais/anexos} % Opcional (anexos)
\end{document}